\section{Results}
\label{sec:results}

\subsection{Headline: data-axis decomposition}

\begin{table}[!t]
\centering
\caption{Macro-F1 on canonical SROIE Task-3 ($n=\VAR{n_test_images}$)
         under three training-set composition regimes.  The marginal
         column reports the F1 lift attributable to each new donor
         corpus added to the joint training pool.  Architecture
         (DONUT) and training recipe (Kim et al.\ 2022) held fixed
         across rows.}
\label{tab:headline}
\begin{tabular}{lcc}
\toprule
Training regime & F1 & marginal $\Delta$ \\ \midrule
SROIE only (single-corpus baseline)         & \VAR{donut_f1_single}        & --- \\
+\,CORD-v2 (schema diversity)               & \VAR{donut_f1_cord}          & \VAR{donut_delta_cord} \\
+\,WildReceipt (imaging diversity)          & \VAR{donut_f1}               & \VAR{donut_delta_wild} \\
\bottomrule
\end{tabular}
\end{table}

\subsection{Per-field decomposition}

\begin{table}[!t]
\centering
\caption{Per-field token F1 on canonical SROIE Task-3 ($n=347$).
         The largest per-field lift attributable to multi-dataset
         training is on \textsf{total} (CORD's 30-field schema
         exercises tax / subtotal / item-line structure that SROIE
         under-represents); the smallest is on \textsf{date}, where
         the regex-conforming substring is already near-saturated
         in the single-corpus baseline.}
\label{tab:perfield}
\begin{tabular}{lcccc}
\toprule
Training regime & Company & Date & Address & Total \\ \midrule
SROIE only            & \VAR{donut_f1_company_single}  & \VAR{donut_f1_date_single}  & \VAR{donut_f1_address_single}  & \VAR{donut_f1_total_single} \\
+\,CORD-v2            & \VAR{donut_f1_company_cord}    & \VAR{donut_f1_date_cord}    & \VAR{donut_f1_address_cord}    & \VAR{donut_f1_total_cord} \\
+\,WildReceipt (full) & \VAR{donut_f1_company}         & \VAR{donut_f1_date}         & \VAR{donut_f1_address}         & \VAR{donut_f1_total} \\
\bottomrule
\end{tabular}
\end{table}

\subsection{Cross-corpus generalisation}

We additionally evaluate the multi-dataset DONUT on the donor
corpora (CORD-v2 test, WildReceipt test) to confirm that the
training-data expansion does not sacrifice SROIE performance for
in-distribution donor performance.

\begin{table}[!t]
\centering
\caption{Cross-corpus macro-F1.  The multi-dataset DONUT remains
         competitive on each donor's own test split, demonstrating
         that the joint training does not collapse onto SROIE at
         the cost of the donor distributions.}
\label{tab:crosscorpus}
\begin{tabular}{lccc}
\toprule
Training regime & SROIE F1 & CORD-v2 F1 & WildReceipt F1 \\ \midrule
SROIE only            & \VAR{donut_f1_single} & \VAR{donut_cord_f1_single}     & \VAR{donut_wild_f1_single} \\
+\,CORD-v2            & \VAR{donut_f1_cord}   & \VAR{donut_cord_f1_cord}       & \VAR{donut_wild_f1_cord} \\
+\,WildReceipt (full) & \VAR{donut_f1}        & \VAR{donut_cord_f1}            & \VAR{donut_wild_f1} \\
\bottomrule
\end{tabular}
\end{table}

\subsection{Training-data manifest provenance}

The multi-dataset training fold consists of \VAR{n_train_total}
receipts: \VAR{n_train_sroie} from SROIE, \VAR{n_train_cord} from
CORD-v2, and \VAR{n_train_wild} from WildReceipt.  Image-level
SHA-256 hashes for the full fold are published in
\texttt{paper1/results/training\_manifest.json}, with per-image
schema mappings that translate each donor corpus's labels into
SROIE's four-field schema (Sec.~\ref{sec:method}).
